\documentclass[11pt]{amsart}
\usepackage[margin=1in]{geometry}
\usepackage{amsmath,amssymb,amsthm,mathtools}
\usepackage{booktabs}
\usepackage{enumitem}
\usepackage[hidelinks]{hyperref}

\newtheorem{theorem}{Theorem}[section]
\newtheorem{lemma}[theorem]{Lemma}
\newtheorem{proposition}[theorem]{Proposition}
\newtheorem{corollary}[theorem]{Corollary}
\theoremstyle{definition}
\newtheorem{definition}[theorem]{Definition}
\newtheorem{remark}[theorem]{Remark}

\DeclareMathOperator{\Tr}{Tr}
\DeclareMathOperator{\Var}{Var}
\DeclareMathOperator{\rk}{rk}
\newcommand{\OK}{\mathcal{O}_K}
\newcommand{\OF}{\mathcal{O}_F}
\newcommand{\Q}{\mathbb{Q}}
\newcommand{\Z}{\mathbb{Z}}
\newcommand{\R}{\mathbb{R}}
\newcommand{\LamK}{\Lambda_K}
\newcommand{\nurel}{\nu_{K/F}}
\newcommand{\file}[1]{\texttt{#1}}

\title{The lifting problem for universal quadratic forms in degrees six and seven}
\date{Draft notes, September 2026. Code and data: \file{higher-degree/} in \url{https://github.com/pav10/2O}}

\begin{document}

\begin{abstract}
We show that no totally real number field of degree $6$ or $7$ admits a universal quadratic
form with rational integer coefficients (a universal $\Z$-form).  Together with the
classification in degrees $\le 5$ by Kala and Yatsyna, the fields of degree at most $7$ with a
universal $\Z$-form are exactly $\Q$, $\Q(\sqrt5)$ and $\Q(\zeta_7+\zeta_7^{-1})$.
The proof replaces the house bound of the degree-$\le 5$ argument by the geometry of the
\emph{variance lattice} $\LamK=\OK/\Z$: a tensor-product minimum bound turns the existence of a
universal $\Z$-form into lower bounds for traces of totally positive integers, and a
``downspread'' argument bounds the shortest vector of $\LamK$ by $(5+\sqrt{21})/2$.
This makes the search space small: in the case where the shortest vector generates $K$,
no number field has to be computed at all (every candidate dies inside the order $\Z[x_1]$);
in the imprimitive case only $18$ fields are ever built, and the last five are excluded by
explicit elements.  A second, independent route through the Mordell--Ko theorem and the
$E_6$ trace threshold confirms every final certificate.  In degree $7$ (no proper subfields,
$\nu(K)\le(11+\sqrt{105})/4$) all $5.3\cdot10^8$ candidate polynomials are excluded inside
$\Z[x_1]$, again without computing a single number field.
\end{abstract}

\maketitle

\section{Statement and overview}\label{sec:intro}

Let $K$ be a totally real number field of degree $d$ with ring of integers $\OK$.  A
\emph{$\Z$-form} is a positive definite quadratic form $Q(x)=\sum_{i\le j}a_{ij}x_ix_j$ with
$a_{ij}\in\Z$; it is \emph{universal over $K$} if it represents every totally positive element
of $\OK$ over $\OK$.  Kala and Yatsyna \cite{KY-EB} proved that for $d\le5$ such a form exists
only for $K=\Q,\Q(\sqrt5),\Q(\zeta_7+\zeta_7^{-1})$.  Their proof uses the Mordell--Ko
theorem \cite{Ko}: in rank $\le5$ every positive definite classical integral form is a sum of
squares of integral linear forms.  This fails in rank $6$ (the lattice $E_6$), which is why
degree $6$ was the first open case.

\begin{theorem}\label{thm:main}
No totally real number field of degree $6$ admits a universal $\Z$-form.
\end{theorem}

\begin{theorem}\label{thm:seven}
No totally real number field of degree $7$ admits a universal $\Z$-form.
\end{theorem}

\begin{theorem}\label{thm:eight}
No totally real number field of degree $8$ admits a universal $\Z$-form.
\end{theorem}

\begin{corollary}\label{cor:le6}
A totally real field of degree $d\le 8$ admits a universal $\Z$-form if and only if it is
$\Q$, $\Q(\sqrt5)$ or $\Q(\zeta_7+\zeta_7^{-1})$.
\end{corollary}

The proof has a theoretical part (\S\S\ref{sec:setup}--\ref{sec:e6}), which reduces the
theorem to a finite computation with explicit bounds, and a computational part
(\S\ref{sec:computation} for degree $6$, \S\ref{sec:deg7} for degree $7$,
\S\S\ref{sec:d8A}--\ref{sec:d8B} for degree $8$, which uses the sharpenings of \S\ref{sec:sharp}).  The main
steps are:
\begin{enumerate}[label=(\roman*)]
\item The necessary condition~(N) of \cite{KY1}: every totally positive $\alpha$ is
  $\alpha=\sum M_{ij}\omega_i\omega_j$ for a positive semidefinite half-integral $M$.
\item A tensor-product bound (Lemma~\ref{lem:tensor}) turns (N) into a \emph{budget--cost}
  inequality $\tau(\alpha)\ge\nu(K)$ (Proposition~\ref{prop:budget}) and into a
  \emph{rigidity} statement for elements with $\tau(\alpha)<\frac32\nu(K)$
  (Proposition~\ref{prop:rigidity}), plus relative versions over subfields.
\item Applying budget--cost to $x_1-n$ and $-x_1-n'$, for $x_1$ a shortest vector of the variance
  lattice, gives $\nu(K)\le(5+\sqrt{21})/2$ (Corollary~\ref{cor:nubound}).  Over a proper
  subfield, the same argument together with a ``positive covering constant'' $\delta_F$ bounds
  the relative minimum (Corollary~\ref{cor:relbound}).
\item The resulting finite enumerations are run in \S\ref{sec:computation}.  Every candidate
  field is excluded, and every exclusion is witnessed by an explicit element.
\end{enumerate}

\section{Setup}\label{sec:setup}

Throughout, $K$ is totally real of degree $d$ with embeddings $\sigma_1,\dots,\sigma_d$.
For $x\in K$ put
\[
\tau(x)=\tfrac1d\Tr_{K/\Q}(x),\qquad
\Var(x)=\tfrac1d\sum_{j=1}^d\bigl(\sigma_j(x)-\tau(x)\bigr)^2 .
\]
Both are intrinsic: if $x$ lies in a subfield $F$, they coincide with the quantities computed
in $F$.  The \emph{variance lattice} is $\LamK=\OK/\Z$ with the positive definite quadratic
form $d\Var$; it is a lattice of rank $d-1$ (the embedding $x\mapsto(\sigma_j(x)-\tau(x))_j$
identifies it with a lattice in $\mathbf 1^\perp\subset\R^d$).  Put
\[
\nu(K)=\min\{\Var(x):x\in\OK\setminus\Z\},\qquad
t(K)=\inf\{\tau(\alpha):\alpha\in\OK^+\setminus\Z\}.
\]
A \emph{shortest vector} is an $x_1\in\OK\setminus\Z$ with $\Var(x_1)=\nu(K)$; it is determined
by its class up to $x_1\mapsto\pm x_1+k$ ($k\in\Z$).

For a subfield $F\subseteq K$ with $f=[F:\Q]$ and $e=[K:F]$, group the embeddings of $K$ into
$f$ blocks $B_1,\dots,B_f$ of size $e$ (the fibres over the embeddings $\rho_1,\dots,\rho_f$
of $F$), and define the block means $\mu_k(x)=\frac1e\sum_{j\in B_k}\sigma_j(x)$ and
\[
\Var_{K/F}(x)=\tfrac1d\sum_{k=1}^f\sum_{j\in B_k}\bigl(\sigma_j(x)-\mu_k(x)\bigr)^2,
\qquad \nurel=\min\{\Var_{K/F}(x):x\in\OK\setminus F\}.
\]
The map $x\mapsto(\sigma_j(x)-\mu_k(x))_{j\in B_k,\,k}$ has kernel $F$, and $\OK/\OF$ is
torsion free, so $\Var_{K/F}$ is a positive definite form on the lattice $\OK/\OF$ of rank
$d-f$.

\section{Condition (N) and the tensor bound}\label{sec:tensor}

\begin{lemma}[{\cite[Prop.~3.1]{KY1}}]\label{lem:N}
Let $\omega_1,\dots,\omega_d$ be a $\Z$-basis of $\OK$.  If $K$ admits a universal $\Z$-form,
then every $\alpha\in\OK^+$ satisfies
\[
\text{(N)}\qquad \alpha=\sum_{i,j}M_{ij}\,\omega_i\omega_j\quad\text{for some positive
semidefinite } M=M^t\text{ with } M_{ii}\in\Z,\ 2M_{ij}\in\Z .
\]
\end{lemma}

\begin{proof}
If $Q$ has Gram matrix $M_Q$ and $\alpha=Q(v)$ with $v\in\OK^r$, write $v=V\omega$ with
$V\in\Z^{r\times d}$ and take $M=V^tM_QV$.
\end{proof}

Equivalently, $2\alpha=Q_{L}(v)$ where $L=(\Z^d,2M)$ is an \emph{even} lattice.  Replacing
$L$ by its quotient by the radical, we may take $L$ positive definite of rank $\le d$; then
$\min L\ge2$, and every vector of $L$ that is not a minimal vector has norm $\ge4$.

\begin{lemma}[tensor floor]\label{lem:tensor}
Let $L,M$ be positive definite lattices (real quadratic forms on free $\Z$-modules), and let
$w\in L\otimes M$, $w\ne0$, have tensor rank $s$.  Then
$Q(w)\ge\frac{s}{\gamma_s^2}\min L\cdot\min M$, where $\gamma_s$ is Hermite's constant.
In particular, for $s\le 8$,
\[
Q(w)\ \ge\ \min L\cdot\min M,\qquad\text{and}\qquad Q(w)\ge\tfrac32\min L\cdot\min M\ \text{ if } s\ge2 .
\]
\end{lemma}

\begin{proof}
By the Smith normal form, the coefficient matrix of $w$ is $\sum_{i=1}^s x_iy_i^t$ with
$x_1,\dots,x_s\in L$ and $y_1,\dots,y_s\in M$ linearly independent.  With the Gram
matrices $X$ of the $x_i$ and $Y$ of the $y_i$ we have $Q(w)=\operatorname{tr}(XY)
=\operatorname{tr}(X^{1/2}YX^{1/2})\ge s(\det X\det Y)^{1/s}$ by the arithmetic--geometric
mean inequality, and $\det X\ge(\min L/\gamma_s)^s$, $\det Y\ge(\min M/\gamma_s)^s$ by
Hermite's inequality for the lattices spanned by the $x_i$ and the $y_i$.  The known values
$\gamma_1^2=1$, $\gamma_2^2=\frac43$, $\gamma_3^2=2^{2/3}$, $\gamma_4^2=2$,
$\gamma_5^2=8^{2/5}$, $\gamma_6^2=(64/3)^{1/3}$, $\gamma_7^2=64^{2/7}$, $\gamma_8^2=4$ (see
\cite{CS}) give $s/\gamma_s^2=1,\ 1.5,\ 1.889,\ 2,\ 2.176,\ 2.163,\ 2.133,\ 2$ for
$s=1,\dots,8$.
\end{proof}

This is the elementary case of Kitaoka's theorem on minima of tensor products \cite{Ki}.
Below, the tensor rank of an element of $L\otimes\LamK$ (or of $L\otimes\OK/\OF$) is at most
$\rk\LamK=d-1$, so Lemma~\ref{lem:tensor} applies whenever $d\le9$.  This is the only place
where the degree enters; all statements of \S\S\ref{sec:tensor}--\ref{sec:bounds} below that
assume $d\le9$ hold for that reason (degrees $6,7,8$ are proved in this note; $d=9$ is
discussed in \S\ref{sec:status}).

\begin{proposition}[budget--cost]\label{prop:budget}
Let $d\le9$ and suppose $K$ admits a universal $\Z$-form.  Then $\tau(\alpha)\ge\nu(K)$ for
every $\alpha\in\OK^+\setminus\Z$; that is, $t(K)\ge\nu(K)$.
\end{proposition}

\begin{proof}
Write $2\alpha=Q_L(v)$ with $v\in L\otimes\OK$ as above, and let $m=\frac1d\sum_j\sigma_j(v)\in
L\otimes\R$.  Averaging $Q_L(\sigma_jv)=Q_L(m)+2B(m,\sigma_jv-m)+Q_L(\sigma_jv-m)$ over $j$,
\begin{equation}\label{eq:decomp}
2\tau(\alpha)=Q_L(m)+\tfrac1d\,Q_{L\otimes\LamK}(\bar v),
\end{equation}
where $\bar v$ is the image of $v$ in $L\otimes\LamK$ and $\LamK$ carries the form $d\Var$.
Since $\alpha\notin\Z$ we have $v\notin L\otimes\Z$, so $\bar v\ne0$ ($\Z$ is saturated in
$\OK$), and Lemma~\ref{lem:tensor} gives $Q_{L\otimes\LamK}(\bar v)\ge 2\cdot d\nu(K)$.
\end{proof}

\begin{proposition}[rigidity]\label{prop:rigidity}
Under the hypotheses of Proposition~\ref{prop:budget}, let $\alpha\in\OK^+\setminus\Z$ with
$\tau(\alpha)<\frac32\nu(K)$.  Then $\alpha=y^2+cy+b$ with $y\in\OK\setminus\Z$, $b,c\in\Z$ and
$\nu(K)\le\Var(y)\le\tau(\alpha)$.
\end{proposition}

\begin{proof}
In \eqref{eq:decomp}, $\frac1dQ(\bar v)<3\nu(K)$.  By Lemma~\ref{lem:tensor}, $\bar v$ has
tensor rank $1$, i.e.\ $\bar v=e\otimes\bar y$, and $Q_L(e)=2$ (otherwise $Q_L(e)\ge4$ and the
cost is $\ge4\nu$).  Lifting, $v=e\otimes y+r$ with $r\in L$, so
$2\alpha=2y^2+2B(e,r)y+Q_L(r)$ with $Q_L(r)\in2\Z$.  Finally $2\tau(\alpha)\ge\frac1dQ(\bar v)=
2\Var(y)$.
\end{proof}

\begin{proposition}[relative versions]\label{prop:relative}
Suppose $K$ ($d\le 9$) admits a universal $\Z$-form, and let $F\subsetneq K$ be a subfield.
Then $\tau(\alpha)\ge\nurel$ for every $\alpha\in\OK^+\setminus F$.  If moreover
$\tau(\alpha)<\frac32\nurel$, then $\alpha=y^2+\gamma y+\beta$ with $y\in\OK\setminus F$,
$\beta,\gamma\in\OF$ and $\Var_{K/F}(y)\le\tau(\alpha)$.
\end{proposition}

\begin{proof}
Decompose $\R^d$ orthogonally into block-constant vectors and their complement.  The part
of $\frac1dQ(\bar v)$ in the complement is $\frac1dQ(\pi\bar v)$, where $\pi\bar v\in
L\otimes\pi(\LamK)$ and $\pi(\LamK)\cong\OK/\OF$ carries the form $d\Var_{K/F}$ of minimum
$d\nurel$.  If $\alpha\notin F$ then $v\notin L\otimes\OF$, so $\pi\bar v\neq0$; apply
Lemma~\ref{lem:tensor} as before.  For rigidity, $\pi\bar v=e\otimes\pi(\bar y)$ with
$Q_L(e)=2$.  Because $\ker\pi$ on $L\otimes\LamK$ is $L\otimes\Lambda_F$, this gives
$v-e\otimes y\in L\otimes\OF$; expand as before, using that $L$ is even.
\end{proof}

\section{Bounding the minima}\label{sec:bounds}

\begin{lemma}[downspread]\label{lem:downspread}
Let $x\in\OK\setminus\Z$ with conjugates $\theta_1<\dots<\theta_d$.  Then
$\alpha_-=x-(\lceil\theta_1\rceil-1)$ and $\alpha_+=-x+(\lfloor\theta_d\rfloor+1)$ are totally
positive and non-rational, and
$\tau(\alpha_-)=\tau(x)-\lceil\theta_1\rceil+1<\tau(x)-\theta_1+1$,
$\tau(\alpha_+)=\lfloor\theta_d\rfloor+1-\tau(x)<\theta_d-\tau(x)+1$.
\end{lemma}

\begin{corollary}\label{cor:nubound}
If $K$ of degree $d\le9$ admits a universal $\Z$-form, then $d\,\nu\ge2(\nu-1)^2$ for
$\nu=\nu(K)$.  In particular $\nu(K)\le\frac{(d+4)+\sqrt{(d+4)^2-16}}{4}$, which for $d=6$
is $(5+\sqrt{21})/2=4.7913\ldots$.
\end{corollary}

\begin{proof}
Let $x_1$ be a shortest vector and $w_j=\sigma_j(x_1)-\tau(x_1)$.  By
Proposition~\ref{prop:budget} and Lemma~\ref{lem:downspread}, $-\min_jw_j>\nu-1$ and
$\max_jw_j>\nu-1$.  Hence $d\nu=\sum w_j^2\ge2(\nu-1)^2$.
\end{proof}

The enumeration of \S\ref{sec:computation} uses the exact forms of these inequalities (with
the integer parts), together with the following consequence of budget--cost for quadratic
polynomials in $x_1$.

\begin{lemma}\label{lem:quadfilter}
Let $x_1$ be a shortest vector of degree $\ge3$ and $P\in\Z[t]$ monic of degree $2$ with
$P(\sigma_j(x_1))>0$ for all $j$.  Under the hypotheses of Proposition~\ref{prop:budget},
$P(\tau(x_1))\ge0$.
\end{lemma}

\begin{proof}
$P(x_1)$ is totally positive and not rational, and $\tau(P(x_1))=\Var(x_1)+P(\tau(x_1))$.
\end{proof}

\begin{definition}
For a totally real field $F$ of degree $f$ with Minkowski lattice $\sigma(\OF)\subset\R^f$,
the \emph{positive covering constant} is
\[
\delta_F=\sup_{u\in\R^f}\ \min\Bigl\{\tfrac1f\textstyle\sum_k(p_k-u_k):\ p\in\sigma(\OF),\
p_k>u_k\ \text{for all }k\Bigr\}.
\]
\end{definition}

For example $\delta_\Q=1$.  The function $h(u)=\min_{p>u}\sum_k(p_k-u_k)$ is
$\sigma(\OF)$-periodic and satisfies $h(u)\le h(u')+\sum_k(u'_k-u_k)$ whenever $u\le u'$
coordinatewise.  Evaluating $h$ on a grid of step $\varepsilon$ that covers a fundamental
parallelepiped therefore gives the rigorous bound $\delta_F\le\max_{\mathrm{grid}}h/f+\varepsilon$
(\file{deltaF.py}).  The computation is in double precision but only ever bounds $h$ from
above: a lattice point $p$ is used for a grid point $u'$ only if $p_k>u'_k+10^{-9}$ in floating
point, which implies $p>u'$ exactly, and $10^{-9}$ is added to the result.  The bounds used in
this note are
\begin{center}
\begin{tabular}{lcccccc}
\toprule
$F$ & $\Q(\sqrt5)$ & $\Q(\sqrt2)$ & $\Q(\zeta_7)^+$ & $\Q(\zeta_9)^+$ & $d_F=725$ & $d_F=1600$\\
\midrule
grid step & $0.0005$ & $0.0005$ & $0.01$ & $0.01$ & $0.05$ & $0.1$\\
rigorous bound & $1.30952$ & $1.70761$ & $1.54132$ & $2.14626$ & $1.90188$ & $2.156694$\\
constant used & $1.3110$ & $1.7091$ & $1.5513$ & $2.1529$ & $1.9144$ & $2.1567$\\
\bottomrule
\end{tabular}
\end{center}
Every constant used is at least the rigorous bound.  (An earlier version took the bounds from
coarser grids rounded to four decimals; at that rounding the first four were about
$2\cdot10^{-5}$ too small, and the coarse grid for $d_F=725$ missed the maximum.  The finer
grids above remove both issues without changing any enumeration.)

\begin{lemma}[relative downspread]\label{lem:reldown}
For $y\in\OK\setminus F$ with block minima $m_k=\min_{j\in B_k}\sigma_j(y)$, there is
$\gamma\in\OF$ with $y+\gamma$ totally positive and
$\tau(y+\gamma)\le\frac1f\sum_k(\mu_k(y)-m_k)+\delta_F$.  The same holds for $-y$ with the
block maxima.
\end{lemma}

\begin{proof}
$\tau(y+\gamma)=\frac1f\sum_k(\mu_k(y)-m_k)+\frac1f\sum_k(\rho_k(\gamma)+m_k)$, and the second
sum can be made $\le f\delta_F$ with all $\rho_k(\gamma)+m_k>0$.
\end{proof}

\begin{corollary}\label{cor:relbound}
If $K$ ($d\le9$) admits a universal $\Z$-form and $F\subsetneq K$ with $e=[K:F]$, then
$e\,\nurel\ge2(\nurel-\delta_F)^2$.
\end{corollary}

\begin{proof}
Let $y$ realise $\nurel$, and write $u_k$ for the vector of within-block deviations in block
$k$, with maximum $M_k$ and minimum $-m'_k$.  By Proposition~\ref{prop:relative} and
Lemma~\ref{lem:reldown}, $D_\pm\ge\nurel-\delta_F$, where $D_+=\frac1f\sum M_k$ and
$D_-=\frac1f\sum m'_k$.  Since $\|u_k\|^2\ge M_k^2+m_k'^2\ge(M_k+m'_k)^2/2$,
$d\nurel=\sum_k\|u_k\|^2\ge\frac1{2f}\bigl(\sum_k(M_k+m'_k)\bigr)^2\ge2f(\nurel-\delta_F)^2$.
\end{proof}

For $F=\Q$ this is Corollary~\ref{cor:nubound} (with $\delta_\Q=1$, $e=d$).

\section{Rigidity via square roots, and the \texorpdfstring{$E_6$}{E6} threshold}\label{sec:e6}

For testing, rigidity is used in the following form, which needs no search in $\OK$.

\begin{lemma}\label{lem:sqshape}
Let $\alpha$ be as in Proposition~\ref{prop:rigidity} and suppose $\Q(\alpha)=K$.  Then there
is $\kappa\in\Z$ with $\kappa\equiv0,3\pmod4$ such that $4\alpha-\kappa=z^2$ with $z\in\OK$,
$4\nu(K)\le\Var(z)\le4\tau(\alpha)$ and
$0<4\min_j\sigma_j(\alpha)-\kappa\le\max\bigl(\frac{d^2}{d-1}\tau(\alpha),\,\Var(\alpha)/\nu(K)\bigr)$
(for $d=6$: $\frac{36}5\tau(\alpha)=7.2\,\tau(\alpha)$).
Equivalently, the polynomial $4^d m_\alpha((t^2+\kappa)/4)\in\Z[t]$ has an irreducible factor
of degree $d$ whose roots have variance in $[4\nu(K),4\tau(\alpha)]$.
\end{lemma}

\begin{proof}
Take $z=2y+c$ and $\kappa=4b-c^2$.  Put $u=4\min\sigma(\alpha)-\kappa$, so that
$|\sigma_j(z)|\ge\sqrt u$.  If $p\ge1$ conjugates of $z$ are positive and $q\ge1$ negative, then
$\Var(z)\ge\frac{pq}{d^2}(2\sqrt u)^2\ge\frac{d-1}{d^2}\,4u$, so $u\le\frac{d^2}{d-1}\tau(\alpha)$.  If they all agree,
$\Var(z)=\Var(\sqrt{4\alpha-\kappa})\le4\Var(\alpha)/u$, and $\Var(z)\ge4\nu$ gives
$u\le\Var(\alpha)/\nu$.
\end{proof}

The second tool is used as an independent cross-check of the final certificates.  It is the
place where $E_6$ enters.

\begin{proposition}[$E_6$ threshold]\label{prop:e6}
Let $d=6$ and $\alpha\in\OK^+$ with $\Tr\alpha<3(3d_K)^{1/6}$.  Then (N) holds for $\alpha$ if
and only if
\[
\text{(SoS$_2$)}\qquad 2\alpha=\sum_k x_k^2\ \text{ with } x_k\in\OK\ \text{ and }\ \sum_kx_k\in2\OK .
\]
\end{proposition}

\begin{proof}
If (SoS$_2$) holds, write $x=U\omega$ and take $M=\frac12U^tU$: its diagonal is integral
because $\sum_kU_{ki}^2\equiv\sum_kU_{ki}\pmod 2$.  Conversely, let $2\alpha=Q_L(v)$ with $L$
even of rank $r\le6$, and write $v=\sum_kf_k\otimes\omega_k$ with $f_k\in L$.  If
$f_1,\dots,f_6$ are linearly dependent, then $2\alpha$ is represented by an even lattice of
rank $\le5$; by Mordell--Ko \cite{Ko,Mo} its Gram matrix is $U^tU$ with $U$ integral, which
gives (SoS$_2$).  Otherwise $G=(B(f_k,f_l))$ is the Gram matrix of an even positive definite
lattice of rank $6$.  Such a lattice has determinant $\ge3$: even unimodular positive
lattices have rank $\equiv0\pmod8$ \cite[Ch.~V]{Se}, and an even lattice of rank $6$ and
determinant $2$ has discriminant form $q\equiv\frac12$ or $\frac32\pmod 2$; gluing it with
$E_7$ (discriminant form $\frac32$), respectively with $A_1$ (discriminant form $\frac12$),
would produce an even unimodular lattice of rank $13$, respectively $7$.  Hence
$2\Tr\alpha=\operatorname{tr}(GT)\ge6(\det G\det T)^{1/6}\ge6(3d_K)^{1/6}$, where $T$ is the
trace form of the basis ($\det T=d_K$).
\end{proof}

The extremal case of the determinant bound is $E_6$ (determinant $3$).  It is the only obstruction to the degree-$\le5$ method in degree $6$, but
it lives above trace $3(3d_K)^{1/6}\ge29.47$.  Every certificate below has trace $\le14$.

\section{The computation}\label{sec:computation}

All code is in \file{higher-degree/src}.  \file{run\_degree6.sh} reproduces everything in
about 10 minutes on 4 cores; \file{run\_sanity.sh} runs the same machinery in degrees $2$--$5$ (gcc, PARI/GP 2.15, python3 with numpy).  Floating point is used
only to \emph{locate} candidates, always with margins in the direction that keeps more
candidates.  Every exclusion is decided in exact arithmetic.

Let $K$ be a sextic field with a universal $\Z$-form and $x_1$ a shortest vector.  Either
$\Q(x_1)=K$ (Case A) or $F=\Q(x_1)$ is a quadratic or cubic subfield (Case B).

\subsection{Case A: the shortest vector generates \texorpdfstring{$K$}{K}}\label{sec:caseA}
Normalise $x_1\mapsto\pm x_1+k$ so that $\Tr x_1\in\{0,1,2,3\}$.  By
Corollary~\ref{cor:nubound}, the minimal polynomial $x^6+a_1x^5+\dots+a_6$ of $x_1$ satisfies
$a_1\in\{0,-1,-2,-3\}$ and $\sum_j(\sigma_jx_1-\tau x_1)^2\le 15+3\sqrt{21}$.
\file{enum6.c} enumerates all real-rooted integral polynomials with these constraints.  It
uses Robinson's method: by Rolle's theorem each coefficient is confined to an interval
determined by the roots of the previous derivative \cite{Ro,MS}.  The program keeps only
those whose roots pass the exact downspread inequalities of Lemma~\ref{lem:downspread} (with
$\nu=\Var(x_1)$) and Lemma~\ref{lem:quadfilter}.

\begin{center}
\begin{tabular}{lr}
\toprule
real-rooted polynomials enumerated (leaves)& 2\,088\,667\\
passing downspread and quadratic filters & 40\,764\\
irreducible & 15\,720\\
\midrule
NS: a shorter non-rational element exists in $\Z[x_1]$ & 873\\
BC: some $\alpha\in\Z[x_1]^+\setminus\Z$ has $\tau(\alpha)<\Var(x_1)$ & 5\,062\\
RG: some $\alpha\in\Z[x_1]^+$ in the rigidity window fails Lemma~\ref{lem:sqshape} & 9\,785\\
\textbf{survivors} & \textbf{0}\\
\bottomrule
\end{tabular}
\end{center}

Each of the three order-level tests is a necessary condition \emph{for $x_1$ to be a shortest
vector of $\LamK$} for a field with a universal $\Z$-form.  They are stated in the order
$\Z[x_1]\subseteq\OK$, so no maximal order is ever computed in Case A.  The list of the
15\,720 polynomials with the reason for each exclusion is \file{results/prim\_status.txt}.

\subsection{Case B: the shortest vector lies in a proper subfield}\label{sec:caseB}
Now $\nu(K)=\nu(F)$, and every element of $F$ is an element of $K$ with the same $\tau$ and
$\Var$.  The same enumeration in degrees $2$ and $3$ (in degree $2$ without
Lemma~\ref{lem:quadfilter}, which needs $\deg x_1\ge3$), using only the NS and BC tests (not
rigidity, since the $y$ of Proposition~\ref{prop:rigidity} may lie outside $F$), leaves
exactly
\[
F\in\{\Q(\sqrt5),\ \Q(\sqrt2),\ \Q(\zeta_7)^+,\ \Q(\zeta_9)^+\},
\qquad \nu(F)=\tfrac54,\ 2,\ \tfrac{14}9,\ 2 .
\]
For each $F$ let $y$ be a relative shortest vector, so $K=F(y)$ since $[K:F]\in\{2,3\}$ is
prime.  Corollary~\ref{cor:relbound} and the certified bounds for $\delta_F$ give:

\begin{center}
\begin{tabular}{lccccc}
\toprule
$F$ & $e$ & $\delta_F\le$ & $\nurel\le$ & candidates & order-level survivors\\
\midrule
$\Q(\sqrt5)$ & 3 & 1.3110 & 3.6513 & 2504 relative cubics & 20\\
$\Q(\sqrt2)$ & 3 & 1.7091 & 4.2272 & 3873 relative cubics & 34\\
$\Q(\zeta_7)^+$ & 2 & 1.5513 & 3.3934 & 1582 values of $\Delta$ & 1\\
$\Q(\zeta_9)^+$ & 2 & 2.1529 & 4.2030 & 2242 values of $\Delta$ & 7\\
\bottomrule
\end{tabular}
\end{center}

Here a relative cubic is $y^3+b_1y^2+b_2y+b_3$ over $\OF$, with $b_1$ modulo $3\OF$ and
$\frac16\Tr_F(\frac23b_1^2-2b_2)=\Var_{K/F}(y)$.  A relative quadratic is
$y=(-b_1+\sqrt\Delta)/2$ with $\Delta\in\OF^+$, $b_1$ modulo $2\OF$, $\Delta\equiv b_1^2\pmod4$ and
$\Var_{K/F}(y)=\Tr_F(\Delta)/12$.  The order-level tests in $\OF[y]$ are the absolute ones
(with $\nu(K)=\nu(F)$) together with the relative budget--cost and ``$y$ is relatively
shortest''.  The $62$ survivors give $18$ distinct fields.  Only these $18$ fields are built
in the whole proof.  On $\OK$, $12$ of them fail relative rigidity
(Proposition~\ref{prop:relative}) and one fails absolute rigidity.  The remaining five fail
(N) at an explicit element:

\begin{center}
\begin{tabular}{llccl}
\toprule
$K$ & $d_K$ & $\Tr\alpha$ & $N\alpha$ & minimal polynomial of the witness $\alpha$\\
\midrule
$\Q(\sqrt5,\zeta_7^+)$ & 300125 & 14 & 29 & $x^6-14x^5+73x^4-182x^3+227x^2-133x+29$\\
$\Q(\sqrt2,\zeta_7^+)$ & 1229312 & 12 & 8 & $x^2-4x+2$ \quad($\alpha=2-\sqrt2$)\\
$\Q(\sqrt2,\zeta_9^+)$ & 3359232 & 12 & 8 & $x^2-4x+2$ \quad($\alpha=2-\sqrt2$)\\
$x^6-9x^4+24x^2-17$ & 7138368 & 12 & 9 & $x^3-6x^2+9x-3$ \quad($\alpha=2-\zeta_9-\zeta_9^{-1}$)\\
$x^6-9x^4+24x^2-19$ & 7978176 & 12 & 9 & $x^3-6x^2+9x-3$ \quad($\alpha=2-\zeta_9-\zeta_9^{-1}$)\\
\bottomrule
\end{tabular}
\end{center}

Each failure of (N) is established by the exhaustive search of \file{condN.gp}.  It uses
$\operatorname{tr}M\le\Tr\alpha/\lambda_{\min}(T)$, the bound
$M_{ii}\le(\sum_j|\sigma_j(\omega_i^*)|\sqrt{\sigma_j\alpha})^2$ with $\omega^*$ the
trace-dual basis, integral solution of the linear system for the off-diagonal entries, a
Fincke--Pohst enumeration of the solution coset, and an exact positive-semidefiniteness
test.  The witnesses are checked a second time by Proposition~\ref{prop:e6}, with a
different search (\file{sos2.gp}).  In the last two fields there is no $x\ne0$ with
$x^2\preceq2\alpha$ at all, and in $\Q(\sqrt5,\zeta_7^+)$ the only such $x$ are $\pm1$.
The first witness coincides with the one found earlier by independent code in the working
pack \cite{Pack}.

\subsection{Proof of Theorem~\ref{thm:main}}
Suppose $K$ is sextic with a universal $\Z$-form and let $x_1$ be a shortest vector.  In
Case A, the minimal polynomial of a suitable normalisation of $x_1$ is among the
$15\,720$ polynomials of \S\ref{sec:caseA}, and each of them violates a condition that a
shortest vector must satisfy.  In Case B, $F=\Q(x_1)$ is one of the four fields listed, and
$K=F(y)$ for a relative shortest vector $y$ within the enumerated range.  The resulting field
is one of the $18$ fields of \S\ref{sec:caseB}, each of which violates rigidity or (N).
\qed

\section{Degree seven}\label{sec:deg7}

Let $K$ be a septic field with a universal $\Z$-form.  Budget--cost, rigidity and
Corollary~\ref{cor:nubound} hold for $d=7$: the tensor rank is at most $\rk\LamK=6$, and
$s/\gamma_s^2\ge\frac32$ for $2\le s\le6$.  Thus $\nu(K)\le(11+\sqrt{105})/4=5.3117\ldots$.
Since $7$ is prime, $K$ has no proper subfields.  So the shortest vector $x_1$ generates $K$
(only Case A occurs), and every non-rational element of $K$ generates $K$.  In particular
Lemma~\ref{lem:sqshape} applies to every totally positive non-rational $\alpha$ in the rigidity
window.

The enumeration of \S\ref{sec:caseA} (with $\Tr x_1\in\{0,1,2,3\}$ and centred
$T_2\le7(11+\sqrt{105})/4$) produces $5.3\cdot10^8$ leaves.  So we also run two of the
order-level tests inside the enumerator (\file{enum6.c}, mode \texttt{strong}):
\begin{itemize}[leftmargin=*]
\item[(NS)] for $y=cx_1^2+bx_1$ ($1\le c\le3$) and $y=x_1^3+ax_1^2+bx_1$, with the integers
  $a,b$ near the real minimiser of $\Var(y)$: if $0<\Var(y)<\Var(x_1)$, then $x_1$ is not a
  shortest vector;
\item[(RIG)] Lemma~\ref{lem:sqshape} for the two downspread elements $\alpha_\pm$ of
  Lemma~\ref{lem:downspread}, whenever $\tau(\alpha_\pm)<\frac32\Var(x_1)$.  As $K=\Q(\alpha)$,
  the conjugates of $z$ are $\pm\sqrt{4\sigma_j(\alpha)-\kappa}$.  So $z\in\OK$ exists iff, for
  some $\kappa$ in the range of the lemma and some choice of signs, all elementary symmetric
  functions of these numbers are integers, and $\Var(z)\in[4\Var(x_1),4\tau(\alpha)]$.
\end{itemize}

\begin{center}
\begin{tabular}{lr}
\toprule
real-rooted polynomials enumerated (leaves) & 533\,562\,406\\
excluded by the downspread inequalities (Lemma~\ref{lem:downspread}) & 531\,897\,850\\
excluded by Lemma~\ref{lem:quadfilter} & 41\,498\\
excluded by (NS) in the enumerator & 2\,532\\
excluded by (RIG) in the enumerator & 1\,619\,004\\
remaining & 1\,522\\
\midrule
irreducible among them & 23\\
exact tests in $\Z[x_1]$ (\S\ref{sec:caseA}): NS / BC / RG & 5 / 3 / 15\\
\textbf{survivors} & \textbf{0}\\
\bottomrule
\end{tabular}
\end{center}

\begin{proof}[Proof of Theorem~\ref{thm:seven}]
A suitable normalisation of the shortest vector $x_1$ of $K$ has its minimal polynomial among
the enumerated leaves.  Every irreducible leaf violates a necessary condition for being the
minimal polynomial of a shortest vector of a field with a universal $\Z$-form.
\end{proof}

No number field and no maximal order is computed in degree $7$.  The in-enumerator tests (NS),
(RIG) are the only new ingredients.  Each exclusion they make on an irreducible polynomial was
re-checked in exact arithmetic (\file{verify\_strong.gp}): for (NS), the exact minimum of $\Var$
on $\Z[x_1]\setminus\Z$ by Fincke--Pohst; for (RIG), \texttt{sqshape}, i.e.\ factorisation of
$4^7m_\alpha((t^2+\kappa)/4)$ over $\Q$.  In degree $7$ this re-check covers
$857\,103$ irreducible polynomials (the other $764\,433$ excluded polynomials are reducible), with no disagreement.  In degree $6$ it covers all
$15\,714$ exclusions on irreducible polynomials made by the strong mode, also with no
disagreement.  \file{run\_degree7.sh} reproduces the computation in about 30 minutes on
4 cores.

\section{Sharpenings}\label{sec:sharp}

This section records improvements found after the proofs above.  They are not needed for
Theorems~\ref{thm:main} and~\ref{thm:seven}.  They shrink the computation considerably, and
they are used for Theorem~\ref{thm:eight} (\S\S\ref{sec:d8A}--\ref{sec:d8B}).

\subsection{A single universal form}
Oh \cite{Oh} determined minimal-rank $n$-universal lattices: $E_6\perp I_7$ and $E_7\perp I_6$
are $6$-universal \cite[Thms.~2.3, 2.5]{Oh}; $E_8\perp I_7$ and $E_7\perp I_8$ are
$7$-universal \cite[Cor.~3.2, Thm.~3.3]{Oh}; $E_8\perp I_8$ is $8$-universal
\cite[Thm.~3.1]{Oh}; and $I_{n+3}$ is $n$-universal for $n\le5$.

\begin{proposition}\label{prop:oh}
Let $d\le8$, and let $U_d$ be $I_8$ ($d\le5$), $E_6\perp I_7$ ($d=6$), $E_8\perp I_7$ ($d=7$),
or $E_8\perp I_8$ ($d=8$).  Let $U_d^{\mathrm{ev}}$ be its even sublattice: $D_8$, $E_6\perp D_7$,
$E_8\perp D_7$, $E_8\perp D_8$ respectively.  For $\alpha\in\OK^+$, condition (N) holds if and
only if $\alpha$ is represented over $\OK$ by the $\Z$-form $\Phi_d=\frac12Q_{U_d^{\mathrm{ev}}}$.  In
particular, $K$ admits a universal $\Z$-form if and only if $\Phi_d$ is universal over $K$.
\end{proposition}

\begin{proof}
If (N) holds, then $2\alpha=Q_L(v)$ with $L$ even, positive definite, of rank $\le d$.  By
$d$-universality there is an isometric embedding $L\to U_d$, and its image consists of vectors
of even norm, hence lies in $U_d^{\mathrm{ev}}$.  The Gram matrix of $\Phi_d$ is half-integral with
integral diagonal, so the converse is Lemma~\ref{lem:N} applied to $Q=\Phi_d$.
\end{proof}

So in degree $6$, besides the sums of squares with parity (the $D_7$ part) only one extra
lattice appears, and it is $E_6$.  This is the structural form of
Proposition~\ref{prop:e6}.  (N) can also be decided by searching representations by one fixed
form of rank $\le 16$.

\subsection{Sharpened rigidity and the square dichotomy}
In the proof of Proposition~\ref{prop:rigidity}, $c=B(e,r)$ and $b=Q_L(r)/2$, so by the
Cauchy--Schwarz inequality in $L$
\[
\kappa:=4b-c^2=Q_L(e)Q_L(r)-B(e,r)^2\ \ge\ 0 .
\]
Moreover $\tau(\alpha)=\Var(y)+(\tau(y)+c/2)^2+\kappa/4$, where $\Var(y)\ge\nu(K)$.

\begin{proposition}\label{prop:sharp}
Under the hypotheses of Proposition~\ref{prop:rigidity} we have $4\alpha=z^2+\kappa$ with
$z\in\OK$, $\kappa\in\Z_{\ge0}$, $\kappa\equiv0,3\pmod4$ and
$\kappa+\tau(z)^2\le4(\tau(\alpha)-\nu(K))$.  Consequently:
\begin{enumerate}[label=(\alph*)]
\item\emph{(square dichotomy)} if $\tau(\alpha)<\min(\frac32\nu(K),\nu(K)+\frac34)$, then
  $\alpha$ is a square in $\OK$;
\item\emph{(relative version)} if $F\subsetneq K$ and $\alpha\notin F$ has
  $\tau(\alpha)<\frac32\nurel$, then $\alpha=y^2+\gamma y+\beta$ as in
  Proposition~\ref{prop:relative}, with $4\beta-\gamma^2$ totally non-negative in $\OF$.
\end{enumerate}
\end{proposition}

\begin{proof}
If $\kappa\ne0$ then $\kappa\ge3$ and $\tau(\alpha)\ge\nu+\frac34$.  If $\kappa=0$, then $c$ is
even and $\alpha=(y+c/2)^2$.  For (b), apply Cauchy--Schwarz at every real place of $F$.
\end{proof}

The effect is visible at once.  In degrees $2$ and $3$, the order-level filters with the
sharpened test leave exactly $\Q(\sqrt5)$ and $\Q(\zeta_7)^+$, so no (N) search is needed:
$2\pm\sqrt2$ and $2-\zeta_9-\zeta_9^{-1}$ are window elements that are not squares.  In the
degree-$6$ Case B, $17$ of the $18$ fields fail the sharpened rigidity (\file{sharp.gp}), and
only $\Q(\sqrt5,\zeta_7^+)$ needs condition (N).

The square dichotomy also constrains Case B before any enumeration.  If $x_1\in F$, then
$\nu(K)=\nu(F)$ and every window element of $F$ must be a square in $K$.  If $[K:F]$ is odd,
it must already be a square in $F$.  If $[K:F]=2$, a non-square $\alpha$ forces $K=F(\sqrt\alpha)$.
In degree $6$ this gives the following.  $F=\Q(\sqrt2)$ is impossible, since it would force
$\Q(\zeta_{16})^+\subset K$.  $F=\Q(\zeta_9)^+$ forces
$K=F(\sqrt{2-\zeta_9-\zeta_9^{-1}})=\Q(\zeta_{36})^+$, which fails the sharpened rigidity.
Only $F=\Q(\sqrt5)$ and $F=\Q(\zeta_7)^+$ need a relative enumeration.

\subsection{Integrality-refined bound for \texorpdfstring{$\nu$}{nu}}
With the integer parts kept, the downspread conditions read $\theta_1\le L=\lfloor\tau-\nu+1\rfloor$ and
$\theta_d\ge U=\lceil\tau+\nu-1\rceil$.  With $a=\tau-L$ and $b=U-\tau$, Cauchy--Schwarz on the
middle conjugates gives $d\nu\ge a^2+b^2+(a-b)^2/(d-2)$.  For each trace class this is a
smaller range (\file{refined\_bound.py}):
\begin{center}
\begin{tabular}{lcc}
\toprule
$d$ & Corollary~\ref{cor:nubound} & refined, $\Tr x_1=0,1,2,\dots$\\
\midrule
6 & $4.791$ & $4.000,\ 4.139,\ 4.222,\ 4.250$\\
7 & $5.312$ & $4.857,\ 4.694,\ 4.204,\ 4.245$\\
8 & $5.828$ & $5.000,\ 5.109,\ 5.188,\ 5.234,\ 5.500$\\
\bottomrule
\end{tabular}
\end{center}

\subsection{Pruning the enumeration}
The same conditions give $\theta_1-\tau\le-a$ and $\theta_d-\tau\ge b$.  Two consequences can be
tested inside the Rolle tree, and they remove only subtrees without admissible leaves:
\begin{itemize}[leftmargin=*]
\item every even centred power sum satisfies $M_k=\sum_j(\theta_j-\tau)^k\ge a^k+b^k$, and $M_4$,
  $M_6$ are fixed at depths $4$, $6$;
\item the middle conjugates lie in $\tau\pm r$ with $r^2=S-a^2-b^2$.  Since the $i$-th root of
  $f^{(d-m)}$ lies in $[\theta_i,\theta_{i+d-m}]$, the roots $2,\dots,m-1$ of $g_m$ lie in that
  window.
\end{itemize}
The enumerator output is unchanged bit for bit.  In degree $6$ the number of leaves drops
from $2\,088\,667$ to $178\,438$.  In degree $7$ it drops from $533\,562\,406$ to $7\,577\,940$,
and the running time from about $25$ minutes to about $1$ minute on $4$ cores.

\subsection{Degree 8, Case A: the rung split}\label{sec:d8A}
Let $x_1$ realise $\nu=\nu(K)$, with conjugates $\theta_1<\dots<\theta_d$, $\tau=\tau(x_1)$,
and put $m=\min(\frac32\nu,\nu+\frac34)$.  Let $n=\lceil\theta_1\rceil-1$ (the \emph{lower rung}),
so $x_1-n\succ0$ and $\tau_-=\tau-n\ge\nu$ by budget--cost.  Similarly
$n'=\lfloor\theta_d\rfloor+1$ and $\tau_+=n'-\tau\ge\nu$.

\begin{lemma}\label{lem:rung}
Exactly one of the following holds.
\begin{enumerate}[label=(\roman*)]
\item[(X)] $\tau_-\ge m$ and $\tau_+\ge m$; then $\theta_1\le\lfloor\tau-m\rfloor+1$ and
  $\theta_d\ge\lceil\tau+m\rceil-1$.
\item[(YL)] $\tau_-<m$; then $x_1-n=y^2$ with $y\in\OK$, $\Var(y)\ge\nu$,
  $\tau(y)^2\le\tau_--\nu$, $\sum_j y_j^2=d\tau_-$ and
  $\sum_j y_j^4=d\,\Var(x_1)+d\tau_-^2$.
\item[(YU)] $\tau_+<m$; then $n'-x_1=y^2$ with the same conditions.
\end{enumerate}
In cases (YL)/(YU) the element $y$ also satisfies the downspread conditions of
Corollary~\ref{cor:nubound} with respect to $\nu$.
\end{lemma}

\begin{proof}
If $\tau_-<m$, Proposition~\ref{prop:sharp}(a) makes $x_1-n$ a square $y^2$.  Then
$\tau_-=\tau(y^2)=\Var(y)+\tau(y)^2$ with $\Var(y)\ge\nu$ (as $y\notin\Z$).  The power sums
follow from $y_j^2=\theta_j-n$ and $\Var(x_1)=\Var(y^2)$.  If $\tau_-\ge m$, then
$n\le\lfloor\tau-m\rfloor$ and $\theta_1\le n+1$.  The upper end is symmetric.  Both
$\tau_\pm<m$ may hold at once; such $x_1$ are found in both (YL) and (YU).  The last
sentence is budget--cost for $y-\lceil y_1\rceil+1$ and $\lfloor y_d\rfloor+1-y$.
\end{proof}

In mode (X) the extreme conjugates are pushed a whole unit further out than the downspread
alone gives, which the tree pruning of the previous subsection turns into a much smaller
tree.  In modes (YL)/(YU) one enumerates $y$ instead of $x_1$.  Here $p_2(y)=d\tau_-$ is fixed,
$p_4(y)$ is fixed (it pins $a_4$ through Newton's identity), and $|p_1(y)|\le d\sqrt{\tau_--\nu}$.
The polynomial of $x_1$ is then computed exactly (\texttt{\_\_int128}) from that of $y$ as
$G(t\mp n)$ with $G(u^2)=\pm f_y(u)f_y(-u)$, and the usual filters are applied to $x_1$.
All of this is in \file{enum6.c}, argument \texttt{X}, \texttt{YL} or \texttt{YU}.

\paragraph{Validation.}  In degrees $6$ and $7$ the split run was compared with the plain
sharpened run on all subtrees.  In degree $8$ it was compared on all $60$ subtrees with
$S\le24$.  Every polynomial found by one run but not the other is either reducible, or
irreducible and certified in exact arithmetic (\file{verify\_split.gp}): a side with
$\tau_\pm<m$ has no square root in $\OK$, or its square root violates budget--cost.  The
latter happens because the split run applies the extra valid test on $y$.
\begin{center}
\begin{tabular}{lrrr}
\toprule
$d$ & only plain: reducible & only plain: certified & only split (all reducible)\\
\midrule
6 & 10 & 1 & yes\\
7 & 27 & 2 & yes\\
8 ($S\le24$) & 99 & 5 & 46\\
\bottomrule
\end{tabular}
\end{center}

\paragraph{Result.}  The complete degree-$8$ Case A run (all $a_1,a_2$ allowed by
$\nu\le3+2\sqrt2$, Corollary~\ref{cor:nubound}) visits $6\,012\,361$ leaves and takes $37$ seconds on $4$ cores.  It produces
$13\,187$ distinct polynomials, of which $82$ are irreducible.  The exact tests in
$\Z[x_1]$ (\file{order\_filter.gp}) exclude $78$ of them: NS $7$,
BC $1$, RG $70$.  The four remaining orders
\[
x^8-8x^6+16x^4-8x^2+1,\ x^8-8x^6+18x^4-12x^2+2,\ x^8-8x^6+17x^4-10x^2+1,\ x^8-8x^6+19x^4-14x^2+1
\]
all fail the sharpened rigidity in the maximal order (\file{fieldtest2}, RG2).
\emph{Hence Case A of degree $8$ is empty.}

\subsection{Degree 8, Case B: the shortest vector lies in a proper subfield}\label{sec:d8B}
Everything in \S\S\ref{sec:tensor}--\ref{sec:bounds} holds for $d=8$ (tensor rank $\le7$,
$s/\gamma_s^2\ge\frac32$), including Proposition~\ref{prop:relative},
Lemma~\ref{lem:reldown} and Corollary~\ref{cor:relbound}.  Let $F=\Q(x_1)\subsetneq K$, so
$[F:\Q]\in\{2,4\}$.

\begin{lemma}\label{lem:subenum}
Let $x_1$ realise $\nu(K)$ and $F=\Q(x_1)$ of degree $f<d$.  Then $\nu(K)=\nu(F)=\Var(x_1)$,
and $x_1$ satisfies the downspread conditions and Lemma~\ref{lem:quadfilter} of degree $f$:
$f\nu\ge2(\nu-1)^2$, both downspreads of $x_1$ are $\ge\nu$, and $P(\tau(x_1))\ge0$ for every
monic $P\in\Z[t]$ of degree $2$ positive at the conjugates of $x_1$.  Moreover no element
of $\Z[x_1]$ or of $\mathcal O_F$ outside $\Z$ has variance $<\nu$, and none in
$\mathcal O_F^+\setminus\Z$ has $\tau<\nu$ (the tests NS, BC).
\end{lemma}

\begin{proof}
$\Var$ and $\tau$ of an element of $F$ are the same whether computed in $F$ or in $K$ (each
conjugate in $F$ occurs $d/f$ times), so $\nu(K)\le\nu(F)\le\Var(x_1)=\nu(K)$.  The
downspread elements $x_1-n$, $n'-x_1$, $P(x_1)$ and the elements in NS and BC lie in $F\subset K$,
so budget--cost in $K$ (Proposition~\ref{prop:budget}) applies to them, and its
consequences are the stated inequalities, computed from the $f$ distinct conjugates.  The
proof of Corollary~\ref{cor:nubound} then runs with the $f$ distinct conjugates.
Rigidity is \emph{not} available in this form, since the $y$ of
Proposition~\ref{prop:rigidity} may lie outside $F$.
\end{proof}

So $x_1$ appears in the degree-$f$ enumeration (\file{run\_enum.sh}, the files
\file{sanity\_d2.txt}, \file{sanity\_d4.txt}), and only the tests NS and BC may be applied
to it there.

\paragraph{Square forcing} (\file{sqforce8.gp}).  If $L$ is a field with $x_1\in L\subseteq K$,
then every window element of $L$, i.e.\ every totally positive $\alpha\in\mathcal O_L\setminus\Z$ with
$\tau(\alpha)<\min(\frac32\nu,\nu+\frac34)$, is a square in $\OK$
(Proposition~\ref{prop:sharp}(a)).  So a window non-square $\alpha$ of $L$ forces
$L(\sqrt\alpha)\subseteq K$.  Iterating this up to degree $8$ gives:
\begin{itemize}[leftmargin=*]
\item \emph{Quartic $F$.}  The degree-$4$ enumeration with the tests valid in a subfield (NS, BC)
  leaves $10$ fields.  For $6$ of them $\nu(\mathcal O_F)$ is smaller than $\Var(x_1)$ for every
  candidate $x_1$ generating $F$, which contradicts Lemma~\ref{lem:subenum}.  For the other $4$ ($d_F=1957,2048,2304,3981$) the forced
  octic field still has a window non-square, which is impossible.
\item \emph{$F=\Q(\sqrt2)$.}  $2-\sqrt2$ forces $\Q(\zeta_{16})^+$.  Its window then forces
  $\Q(\zeta_{32})^+$, which has a window non-square: impossible.
\item \emph{$F=\Q(\sqrt5)$.}  $\nu=\frac54$ and the only window element is
  $\varphi^2$, a square, so nothing is forced.
\end{itemize}

\paragraph{$F=\Q(\sqrt5)$} (\file{rel8.gp}).  Let $\beta\in\OK\setminus F$ realise
$\nurel$.  Corollary~\ref{cor:relbound} with $e=4$ and $\delta_F\le1.3110$ gives
$\nurel\le4.2142$.  Either $F(\beta)=K$ or $F'=F(\beta)$ is a quartic field.
\begin{itemize}[leftmargin=*]
\item \emph{$F'$ quartic.}  $\beta=(-b+\sqrt D)/2$ with $D\in\OF^+$ and $\Tr D\le8\nurel$.  With
  the exact relative downspread test (Lemma~\ref{lem:reldown} with the actual covering
  value $\delta(m)$ for the target in place of $\delta_F$), only two
  fields remain: $F'$ of discriminant $725$ ($x^4-x^3-3x^2+x+1$, the totally real quartic
  field of least discriminant) and $F'=\Q(\sqrt2,\sqrt5)$ ($d_{F'}=1600$).  In both, every window element is a square,
  so nothing is forced.  $K/F'$ is then a relative quadratic, with $x_1\in F\subset F'$ and
  $\nu_{K/F'}$ bounded by Corollary~\ref{cor:relbound} ($e=2$).  Here
  $\delta_{F'}\le1.9144$ resp.\ $2.1567$ (\file{deltaF.py}, see the table in \S\ref{sec:bounds}), so
  $\nu_{K/F'}\le3.886$ resp.\ $4.21$.  Before the order-level tests we apply the exact
  relative downspread test: no $\gamma\in\mathcal O_{F'}$ may give $\beta'+\gamma\succ0$ with
  $\tau(\beta'+\gamma)<\nu_{K/F'}$, for $\beta'=\pm\beta$ (\texttt{enum\_e2g8}).
\item \emph{$K=F(\beta)$.}  Relative quartics $x^4+c_1x^3+\dots+c_4$ over $\OF$, with
  $c_1$ reduced modulo $4\OF$ and up to the conjugation $\sqrt5\mapsto-\sqrt5$, which leaves
  $10$ residues.  $c_2$ ranges over a box given by the relative variance, and $c_3,c_4$ over
  Rolle intervals at both real places of $F$.  The exact relative downspread test comes
  first, then irreducibility over $F$, and then the relative order-level tests NSR, BCR and
  RGR in $\OF[\beta]$ (\file{rel6.gp}).
\end{itemize}
\paragraph{Result for $F=\Q(\sqrt5)$.}
\begin{itemize}[leftmargin=*]
\item Relative quartics: $468\,986$ polynomials in the boxes.  The exact relative downspread
  test removes $450\,708$ of them.  Of the irreducible ones, the relative order-level tests
  exclude $7\,316$ (NSR $1\,577$, BCR $5\,739$), and $815$ remain.  Running time about
  $15$--$30$ minutes per residue on one core (\file{run\_q4\_par.sh}).
\item $K/F'$ quadratic: $F'$ of discriminant $725$ gives $23\,613$ values of $D$, with $930$
  removed by the relative downspread test, $161$ by NSR/BCR, and $52$ left.  $F'=\Q(\sqrt2,\sqrt5)$
  gives $20\,976$ values of $D$, with $859$ / $108$ / $97$.
\item The $964$ order-level survivors give $390$ distinct octic fields
  (\file{collect8.gp}; list in \file{fields\_all.txt}).  In the maximal order,
  $386$ fail the relative sharpened rigidity over $\Q(\sqrt5)$ (RGR2), and $2$ fail the relative
  BC test.  Two fields survive:
  \begin{itemize}
  \item $K_1$: $x^8-2x^7-12x^6+26x^5+17x^4-36x^3-5x^2+11x-1$, $d_{K_1}=5^4\cdot29^4$.  This is the
    Galois closure (dihedral of order $8$) of the quartic field of discriminant $725$.
  \item $K_2$: $x^8-3x^7-4x^6+13x^5+5x^4-13x^3-4x^2+3x+1$, $d_{K_2}=5^4\cdot29^2\cdot1249$.
    This is a non-Galois quadratic extension of the same quartic field.
  \end{itemize}
  Both fail (N) (\file{condN.gp}) at an element $\alpha$ of trace $18$ in a quartic
  subfield of discriminant $725$, with minimal polynomial $x^4-9x^3+27x^2-31x+11$.  For $K_2$
  it is the $11$th element tested by increasing trace.  For $K_1$ the search by increasing
  trace finds it as the $21$st element.  Six earlier elements of trace $18$ exhaust the
  $2$\,GB stack and are left undecided, which does not matter for the conclusion
  (\file{witness\_k1\_d8.gp}).
\end{itemize}
Together with Case A (\S\ref{sec:d8A}) this proves Theorem~\ref{thm:eight}.  The
floating-point caveats of \S\ref{sec:status} apply, now including $\delta_{F'}$ for the two
quartic fields $F'$ (grid step $0.1$).

\section{Consistency checks and status}\label{sec:status}

\begin{itemize}[leftmargin=*]
\item \textbf{Enumerator.}  \file{enum6.c} agrees exactly with an exact-arithmetic
  re-implementation (\file{enum\_check.gp}, real-rootedness by Sturm sequences) on all 56
  subtrees and all $2\,088\,667$ leaves in degree $6$, and with brute force over coefficient
  boxes in degrees $3$--$5$.  In degree $7$ it agrees with \file{enum\_check.gp} on
  the $48$ subtrees with the fewest leaves ($2\,267\,434$ leaves in total).
\item \textbf{Degrees $\le5$ reproduced.}  The same pipeline, with the degree-$d$ bound of
  Corollary~\ref{cor:nubound}, leaves $\Q(\sqrt2)$, $\Q(\sqrt5)$ in degree $2$, and
  $\Q(\zeta_9)^+$, $\Q(\zeta_7)^+$ in degree $3$.  Of these, $\Q(\sqrt2)$ and
  $\Q(\zeta_9)^+$ fail (N) at $2-\sqrt2$ and $2-\zeta_9-\zeta_9^{-1}$.  Nothing is left in
  degree $5$.  In degree $4$, Case A leaves nothing, and Case B leaves three fields
  ($d_K=725,1600,7168$), all failing (N).  This recovers \cite[Thm.~1.2]{KY-EB}.
\item \textbf{Controls.}  (N) holds for all tested elements of $\Q(\sqrt5)$ and
  $\Q(\zeta_7)^+$; both fields survive every filter.
\item \textbf{Degree 8} (\file{validate8.sh}).
  \begin{enumerate}[label=(\alph*)]
  \item Every exclusion made inside the split enumerator on an irreducible polynomial is
    re-checked in exact arithmetic (\file{verify\_split.gp}, \texttt{vsharp}).  There are
    $108\,159$ of them (NS $24$, sharpened RIG $107\,539$, $y$-budget $596$), with no
    disagreement.  The other $484\,381$ excluded polynomials are reducible.
  \item At every one of the $392$ maximal-order exclusions (Case A: $4$ RG2; Case B:
    $386$ RGR2, $2$ BCR), condition (N) is decided directly at the element named by the test
    (\file{certN8.gp}).  It fails in all $392$ cases, so every excluded octic field has two
    independent certificates.
  \item The relative-quartic candidates of \texttt{enum\_q4} coincide, for all $16$ residues
    of $c_1$, with an independent brute-force search (\file{check\_q4.gp}).  That search uses
    power-sum boxes for $c_2,c_3,c_4$ and an exact real-rootedness test via the norm
    polynomial and Sturm sequences.  No candidate is missing, and none is superfluous.
  \item The rung split was compared with the plain enumerator on all $115$ subtrees of
    Case A (\file{compare\_split.sh}).  Every polynomial found by only one of the runs is
    reducible, or certified exactly: its forced side has no square root, or the square root
    violates budget--cost.  In total $144$ polynomials were found only by the plain run
    ($137$ reducible, $7$ certified), and $46$ only by the split run (all reducible).  The
    slowest plain runs take up to $52$ minutes; the subtree $(0,-20)$ was split $4$ ways by
    $a_3\bmod 4$ (\texttt{ENUM\_A3MOD}).
  \end{enumerate}
  \textbf{Reproduction.}  \file{run\_degree8.sh} was rerun from scratch in a clean copy of the
  repository.  It reproduces every intermediate count and the final field list exactly (fields,
  discriminants and verdicts of all $390$ Case B fields); for $3$ fields a different subfield
  is named in the RGR2 certificate, as the order of \texttt{nfsubfields} varies.
\item \textbf{Floating point.}  Used for embeddings (60 digits), for root brackets in the
  enumerator (80-bit, with margins), and for $\delta_F$ (double precision with a grid margin
  of $0.002$ resp.\ $0.02$).  The grid values $\delta_{\Q(\sqrt5)}\approx1.3090=\varphi^2/2$ and
  $\delta_{\Q(\sqrt2)}\approx1.7071=1+1/\sqrt2$ are consistent with exact values.
\item \textbf{Recommended before publication.}  An independent re-implementation of Case A
  and of the (N) searches (for example in Magma), and an exact-arithmetic computation of
  $\delta_F$ for the four fields.
\end{itemize}

\begin{remark}[degree $9$]
Lemma~\ref{lem:tensor} covers tensor rank $\le8$, so \S\S\ref{sec:tensor}--\ref{sec:bounds},
the sharpened rigidity and the rung split hold for $d=9$, with
$\nu(K)\le(13+\sqrt{153})/4\approx6.34$ (refined per trace class: $\nu\le6$, \file{refined\_bound.py}).
\begin{itemize}[leftmargin=*]
\item \emph{Case B.}  The proper subfields are cubic, $[K:F]=3$ is odd, and a window
  non-square $\alpha$ of $F$ would give $F(\sqrt\alpha)\subseteq K$ of even degree: impossible.  The
  degree-$3$ enumeration with NS/BC leaves $\Q(\zeta_9)^+$, excluded this way by
  $2-\zeta_9-\zeta_9^{-1}$, and $F=\Q(\zeta_7)^+$ ($\nu=14/9$), which needs the relative cubics
  over $F$ with $\nu_{K/F}\le4.0011$ (Corollary~\ref{cor:relbound}, $e=3$, $\delta_F\le1.5513$;
  \file{rel9.gp}, about $3$ hours on $4$ cores).
\item \emph{Case A.}  The split enumerator agrees with the plain one on the $60$ degree-$9$
  subtrees with $S\le24$ (all differences reducible or certified).  But the tree grows too
  fast: the subtree $(0,-12)$ ($S=24$) has $3.3\cdot10^5$ leaves in mode X ($20$ s), $(0,-14)$
  does not finish in $2$ minutes, and $S$ runs up to $54$.  About $90\%$ of the leaves fail only
  the final rung condition, which the Rolle tree cannot test earlier, and pruning with $M_8$
  changes almost nothing.  Degree $9$ therefore needs a further theoretical restriction in
  Case A.
\end{itemize}
\end{remark}

\begin{thebibliography}{KY2}

\bibitem[CS]{CS}
J.~H. Conway, N.~J.~A. Sloane,
\emph{Sphere Packings, Lattices and Groups}, 3rd ed., Grundlehren Math. Wiss. 290,
Springer, 1999.

\bibitem[Ki]{Ki}
Y.~Kitaoka, \emph{Arithmetic of Quadratic Forms}, Cambridge Tracts in Math. 106,
Cambridge University Press, 1993.

\bibitem[Ko]{Ko}
C.~Ko, On the representation of a quadratic form as a sum of squares of linear forms,
Q. J. Math. 1 (1937), 81--98.

\bibitem[KY1]{KY1}
V.~Kala, P.~Yatsyna, Lifting problem for universal quadratic forms,
Adv. Math. 377 (2021), 107497.

\bibitem[KY2]{KY-EB}
V.~Kala, P.~Yatsyna, Even better sums of squares over quintic and cyclotomic fields,
arXiv:2402.03850v2.

\bibitem[MS]{MS}
J.~McKee, C.~Smyth, Salem numbers of trace $-2$ and traces of totally positive algebraic
integers, in \emph{Algorithmic Number Theory}, Lecture Notes in Comput. Sci. 3076,
Springer, 2004, 327--337.

\bibitem[Mo]{Mo}
L.~J. Mordell, The representation of a definite quadratic form as a sum of two others,
Ann. of Math. 38 (1937), 751--757.

\bibitem[Oh]{Oh}
B.-K. Oh, Universal $\Z$-lattices of minimal rank, Proc. Amer. Math. Soc. 128 (2000),
683--689.

\bibitem[Pack]{Pack}
\emph{Flatness, variance, and the lifting problem}, working notes (unpublished).

\bibitem[Ro]{Ro}
R.~Robinson, Algebraic equations with span less than 4, Math. Comp. 18 (1964), 547--559.

\bibitem[Se]{Se}
J.-P. Serre, \emph{A Course in Arithmetic}, Grad. Texts in Math. 7, Springer, 1973.

\end{thebibliography}

\end{document}
